\documentclass[11pt]{article}
\usepackage[margin=1.1in]{geometry}
\usepackage[colorlinks=true,urlcolor=blue]{hyperref}
\usepackage{url}
\setlength{\parindent}{0pt}
\setlength{\parskip}{4pt}
\begin{document}
\begin{center}
{\Large\bf Response to Reviewers (Round 2)}\\[4pt]
{\it Can AI be Easy? Lessons Learned from the EZR.py Toolkit}\\[2pt]
Menzies, Srinivasan, Ganguly \hfill Software: Practice and Experience
\end{center}
\medskip


We thank the reviewer for a second careful reading. All new or
altered text in this revision is again printed in blue. Point by
point:

\subsection*{Reviewer 1}

\smallskip\noindent{\bf M1(a).} {\em The abstract's text-mining claim is
contradicted by \S7.4.} Fixed. The abstract now states that the
recall match holds on three of the four corpora and that on the
fourth (Kitchenham) the filter converges faster but plateaus
below FASTREAD's ceiling, pointing to the fuller treatment in
\S7.

\smallskip\noindent{\bf M1(b).} {\em ``Tested on the 120+ tasks'' overstates
the evidence for results (a) and (c).} Fixed. The abstract now
scopes each result: (a) 20 randomly selected MOOT benchmarks,
(b) the MOOT tasks, (c) Yu et al.'s four SLR corpora. The
matching overstatement in \S8.1 is likewise rescoped: the
124-task claim now attaches only to the active-learning results,
and the optimizer and text-mining studies are scoped explicitly.

\smallskip\noindent{\bf M1(c).} {\em Tables 2 and 4 still present
companion-derived numbers as findings of this paper.} Fixed. The
\S6 row of Table~2 now opens ``A companion
study\ldots (not a result of this paper) reports\ldots'', and the
label-hunger row of Table~4 marks its comparator
numbers as drawn from companion studies.

\smallskip\noindent{\bf M2.} {\em ``Co-authored by the first author''
understates the overlap now that Ganguly is an author.} Fixed.
All instances (abstract, \S6.4.4, \S8.2, \S8.3, \S8.4, and
the closing remarks) now read ``co-authored by authors of this
paper.''

\smallskip\noindent{\bf M3(a).} {\em Cite the version-specific DOI and git
tag.} Fixed, with an apology: DOI
\textsf{10.5281/zenodo.17118009} was in fact already
version-specific (it resolves to the \textsf{v0.5.0} deposit);
our text mislabeled it a concept DOI. The paper now names the
tag and version-specific DOI at every point of use
(\textsf{v0.9.4}, \textsf{10.5281/zenodo.22554751}, freshly
archived for this revision; listings trace to \textsf{v0.5.0},
\textsf{10.5281/zenodo.17118009}) and gives the true concept
DOI (\textsf{10.5281/zenodo.11183058}) for the full release
history.

\smallskip\noindent{\bf M3(b).} {\em Confirm the Zenodo deposit contains the
400-line single-file version.} Confirmed. We downloaded the
\textsf{v0.5.0} deposit from Zenodo
(\textsf{timm/ezr-v0.5.0.zip} under DOI
\textsf{10.5281/zenodo.17118009}) and verified it contains
\textsf{ezr/ezr.py}: 412 physical lines, 44 functions, matching
the listings of this paper. The current head of the repository
(\textsf{v0.9.4}, two files, less than 1000 lines) is now also
archived, under DOI \textsf{10.5281/zenodo.22554751}, and the
paper's size claims now say ``less than 1000 lines'' so that the
text matches the current release. Additionally: a new Appendix~A gives tested
two-minute instructions (\textsf{pip install ezr} plus the MOOT
data) for running the toolkit's fast demos and its full self-test
suite, with sample output, so that an editor or reviewer can
check the code runs without setting up our full experiments.

\smallskip\noindent{\bf M4.} {\em Table 4's ``weeks'' contradicts \S2.6's
``one session,'' undermining the order-of-magnitude claim.}
Fixed. \S2.6 now says ``weeks of sessions vs.\ a
semester course,'' and the order-of-magnitude claim explicitly
excludes the pedagogy axis, noting that the gap there is weeks
vs.\ a semester and rests on author experience rather than
formal evaluation.

\smallskip\noindent{\bf M5(a).} {\em \S8.4's surrogate-bias paragraph is
stale.} Fixed. The paragraph now names panel~(a), corrects the
count (\textsf{sa-r} at 100 on 13 of 20 rows, mean 98, under the
new 100-repeat runs), reports
the panel~(b) view (no cell at 100, mean 94), and downgrades
``decisive evidence'' to ``suggestive'': the surrogate drives
the search to within a few points of the optimum, but panel~(b)
shows the optimum is approached, not exactly attained.

\smallskip\noindent{\bf M5(b).} {\em The regenerated panel (a) shifts several
column means; 20 repeats do not stabilise \textsf{ls-r}.} The
caption of Table~5 now states that these runs
supersede the original submission's table and that high-variance
columns shift several points between regenerations. Further, we
took the first of the reviewer's two offered remedies: the whole
table (all four treatments, both panels, one set of runs) was
regenerated with the repeat count raised from 20 to 100, so each
cell mean now carries a standard error of sd/10. Under the
100-repeat runs the qualitative conclusions are unchanged
(\textsf{ls-r} remains worst and most erratic;
\textsf{ls+r}, mean 97, still does not pass SA, mean 98;
restart on SA remains a rounding error), and all numbers quoted
in \S5.2 and \S8.4 have been updated to the new runs.

\smallskip\noindent{\bf Minor 1.} {\em ``?'' bibliography entries; SHAP and
ReliefF uncited.} Fixed. Five missing database entries
(Lundberg \& Lee 2017 for SHAP; Robnik-\v{S}ikonja \&
Kononenko 1997 for ReliefF; Yu et al.\ 2018 ``Finding better
active learners''; PEP~544; PEP~557) have been added, repairing
all eight ``?'' instances.

\smallskip\noindent{\bf Minor 2.} {\em Broken cross-reference and unbalanced
parenthesis in \S6.2.} Fixed; the reference now resolves to
Figure~1 and the parenthesis is balanced.

\smallskip\noindent{\bf Minor 3.} {\em The \S7 numbers do not reconcile (warm
start 24, budget 50, x-axes to 300--600, plateaus ``after 100
labels,'' 40 labels/108 papers, ``100 vs.\ 300+'').} Fixed; the
reviewer caught real errors, and we thank them. The true protocol
(now stated in \S7.3, and readable from the ``40'' and
``Baseline'' marks on Figure~5 itself) is: warm start
of 40 labeled samples (20 known-relevant plus 20 presumptive
negatives), not 24; then one label per acquisition, with each run
continuing to the number of labels Yu et al.\ spent on that
corpus (290--630), which is what the x-axis counts. The ``budget
was 50'' sentence was a copy-slip from \S6 and has been deleted.
The Kitchenham claim is now stated in the same currency (after
the 40-label warm start, recall is 87\%, settling to 82\% = 108
of 132). Yu et al.'s labeling cost is now given as 300--630 in
both \S7.2 and the \S7.4 summary.

\smallskip\noindent{\bf Minor 4.} {\em Confirmed as fixed.} We
thank the reviewer; no further action was needed on this point.

\smallskip\noindent{\bf Minor 5.} {\em Three identical affiliations.} Fixed;
all three authors now share one affiliation.


\end{document}
